\section{Open Research Directions}

\subsection{Maximizing offline + online handoff}

Varios de los expositores mencionaron hybrid strategies: primero offline, después un fine-tuning online limitado. La pregunta clave es cómo hacer el handoff de manera segura:

\begin{itemize}
    \item Establecer un clear boundary: ¿en qué situaciones se permite query environment?
    \item Evaluar el online potential ya en la offline phase (e.g., Q variance in shifted states).
    \item Usar meta-RL para entrenar la supervisor policy decide when to go online.
\end{itemize}

\subsection{Counterfactual auditing}

Actualmente el governance todavía depende de los replay logs; en el futuro podría adoptarse counterfactual referents: generar hypothetical trajectories to test safety constraints before deployment.

\begin{knowledgebox}{Counterfactual auditing}
Usar un learnt world model para generar la comparación de «qué pasaría si la policy hubiera tomado una acción distinta», lo cual permite que el safety team conozca la failure trajectory potencial y así identificar más rápido el policy drift.
\end{knowledgebox}

\subsection{Resumen de la sección}

Las direcciones de research incluyen el offline→online handoff, counterfactual auditing y una interpretabilidad más fuerte (como los gradient-level tracebacks).

